% GENERATED FILE. Edit canonical lessons, then run npm run generate:books.
% canonical-source-hash: fnv1a64:b1c64b9ce51b664c
% canonical-lessons: KA-C470-vistarisu, KA-C470-pravesadvara, KA-C470-kadita, KA-C470-sarabaraju, KA-C470-duravani

\chapter{City Works and Supplies}
\label{ch:ka-city-works-and-supplies}
\hlchaptermodality{470}

\begin{chapteropening}
\textbf{By the end of this chapter:} \emph{I can say to widen, an entrance gate, a power cut, supplies and a telephone.}

Complete the last lesson of chapter 470: I can say to widen, an entrance gate, a power cut, supplies and a telephone.
\end{chapteropening}

\section[\texorpdfstring{\kn{ವಿಸ್ತರಿಸು} (vistarisu)}{ವಿಸ್ತರಿಸು (vistarisu)}]{\kn{ವಿಸ್ತರಿಸು} (vistarisu) — to widen; to expand}

\label{lesson:KA-C470-vistarisu}



\begin{quote}
Before the new one: say the Kannada for a bakery, then the Kannada for a home delivery.
\end{quote}

\subsection*{You'll want to know: \kn{ವಿಸ್ತರಿಸು}}
\textbf{\kn{ವಿಸ್ತರಿಸು}} — \emph{vistarisu} — \textquotedblleft{}to widen; to expand\textquotedblright{}.

\begin{quote}
\kn{ಪುರಸಭೆ} \kn{ನಮ್ಮ} \kn{ಬೀದಿಯನ್ನು} \kn{ವಿಸ್ತರಿಸಿತು}.
\end{quote}

\begin{quote}
\emph{purasabhe namma bīdiyannu vistarisitu.}
\end{quote}

\begin{quote}
The town council widened our street.
\end{quote}

Past (I …): \kn{ವಿಸ್ತರಿಸಿದೆ}, \emph{vistariside}.

Of roads, shops and towns; \kn{ವಿಸ್ತಾರ} is extent or breadth.

\begin{cousinweb}
\textbf{\kn{ವಿಸ್ತರಿಸು}} (\emph{vistarisu}) is built from Sanskrit \emph{vistara}, \textquotedblleft{}spreading, extension, expansion\textquotedblright{} (Monier-Williams), with the verb ending -\emph{isu}.
\end{cousinweb}

\begin{grammarlens}[title={What you've built}]
One more word for getting around.
\end{grammarlens}

\subsection*{Your turn}
\begin{itemize}
\raggedright
  \item \emph{Say it:} \emph{vistarisu}
  \item \emph{Say it:} \emph{vistarisu}, once more
  \item \emph{Say it:} \emph{purasabhe namma bīdiyannu vistarisitu.} — The town council widened our street.
  \item \emph{Recall:} say the Kannada for a bakery, then the Kannada for a home delivery, then say \emph{vistarisu} again
\end{itemize}

\begin{tcolorbox}[breakable,colback=teal!4,colframe=teal!35!black,title={Before you move on}]
What does \kn{ವಿಸ್ತರಿಸು} mean? (\textquotedblleft{}To widen; to expand\textquotedblright{}.) Say it once more.
\end{tcolorbox}
\section[\texorpdfstring{\kn{ಪ್ರವೇಶದ್ವಾರ} (pravēśadvāra)}{ಪ್ರವೇಶದ್ವಾರ (pravēśadvāra)}]{\kn{ಪ್ರವೇಶದ್ವಾರ} (pravēśadvāra) — an entrance gate}

\label{lesson:KA-C470-pravesadvara}



\begin{quote}
Before the new one: say the Kannada for a home delivery, then the Kannada for to widen.
\end{quote}

\subsection*{You'll want to know: \kn{ಪ್ರವೇಶದ್ವಾರ}}
\textbf{\kn{ಪ್ರವೇಶದ್ವಾರ}} — \emph{pravēśadvāra} — \textquotedblleft{}an entrance gate\textquotedblright{}.

\begin{quote}
\kn{ವಸ್ತುಸಂಗ್ರಹಾಲಯದ} \kn{ಪ್ರವೇಶದ್ವಾರದಲ್ಲಿ} \kn{ಟಿಕೆಟ್} \kn{ಕೊಡುತ್ತಾರೆ}.
\end{quote}

\begin{quote}
\emph{vastusaṅgrahālayada pravēśadvāradalli ṭikeṭ koḍuttāre.}
\end{quote}

\begin{quote}
Tickets are given out at the entrance gate of the museum.
\end{quote}

\kn{ಪ್ರವೇಶ}, entry, plus \kn{ದ್ವಾರ}, a door or gate.

\begin{cousinweb}
\textbf{\kn{ಪ್ರವೇಶದ್ವಾರ}} (\emph{pravēśadvāra}) is built from two Sanskrit words: \emph{pravēśa}, \textquotedblleft{}an entrance\textquotedblright{} (CDIAL 8809), and \emph{dvāra}, \textquotedblleft{}a door\textquotedblright{} (CDIAL 6663).
\end{cousinweb}

\begin{grammarlens}[title={What you've built}]
One more word for getting around.
\end{grammarlens}

\subsection*{Your turn}
\begin{itemize}
\raggedright
  \item \emph{Say it:} \emph{pravēśadvāra}
  \item \emph{Say it:} \emph{pravēśadvāra}, once more
  \item \emph{Say it:} \emph{vastusaṅgrahālayada pravēśadvāradalli ṭikeṭ koḍuttāre.} — Tickets are given out at the entrance gate of the museum.
  \item \emph{Recall:} say the Kannada for a home delivery, then the Kannada for to widen, then say \emph{pravēśadvāra} again
\end{itemize}

\begin{tcolorbox}[breakable,colback=teal!4,colframe=teal!35!black,title={Before you move on}]
What does \kn{ಪ್ರವೇಶದ್ವಾರ} mean? (\textquotedblleft{}An entrance gate\textquotedblright{}.) Say it once more.
\end{tcolorbox}
\section[\texorpdfstring{\kn{ಕಡಿತ} (kaḍita)}{ಕಡಿತ (kaḍita)}]{\kn{ಕಡಿತ} (kaḍita) — a power cut; a cut}

\label{lesson:KA-C470-kadita}



\begin{quote}
Before the new one: say the Kannada for to widen, then the Kannada for an entrance gate.
\end{quote}

\subsection*{You'll want to know: \kn{ಕಡಿತ}}
\textbf{\kn{ಕಡಿತ}} — \emph{kaḍita} — \textquotedblleft{}a power cut; a cut\textquotedblright{}.

\begin{quote}
\kn{ಬೇಸಿಗೆಯಲ್ಲಿ} \kn{ಪ್ರತಿದಿನ} \kn{ಎರಡು} \kn{ಗಂಟೆ} \kn{ಕರೆಂಟ್} \kn{ಕಡಿತ} \kn{ಇರುತ್ತದೆ}.
\end{quote}

\begin{quote}
\emph{bēsigeyalli pratidina eraḍu gaṇṭe kareṇṭ kaḍita iruttade.}
\end{quote}

\begin{quote}
In summer there is a two-hour power cut every day.
\end{quote}

\kn{ಕರೆಂಟ್} \kn{ಕಡಿತ}, a power cut; \kn{ಸಂಬಳ} \kn{ಕಡಿತ}, a pay cut.

\begin{cousinweb}
\textbf{\kn{ಕಡಿತ}} (\emph{kaḍita}) is native Dravidian: DEDR 1125 lists Kannada \emph{kaḍi}, \textquotedblleft{}to cut, to chop\textquotedblright{}, with the noun \emph{kaḍita}, \textquotedblleft{}a cut, a portion cut off\textquotedblright{}.
\end{cousinweb}

\begin{grammarlens}[title={What you've built}]
One more word for getting around.
\end{grammarlens}

\subsection*{Your turn}
\begin{itemize}
\raggedright
  \item \emph{Say it:} \emph{kaḍita}
  \item \emph{Say it:} \emph{kaḍita}, once more
  \item \emph{Say it:} \emph{bēsigeyalli pratidina eraḍu gaṇṭe kareṇṭ kaḍita iruttade.} — In summer there is a two-hour power cut every day.
  \item \emph{Recall:} say the Kannada for to widen, then the Kannada for an entrance gate, then say \emph{kaḍita} again
\end{itemize}

\begin{tcolorbox}[breakable,colback=teal!4,colframe=teal!35!black,title={Before you move on}]
What does \kn{ಕಡಿತ} mean? (\textquotedblleft{}A power cut; a cut\textquotedblright{}.) Say it once more.
\end{tcolorbox}
\section[\texorpdfstring{\kn{ಸರಬರಾಜು} (sarabarāju)}{ಸರಬರಾಜು (sarabarāju)}]{\kn{ಸರಬರಾಜು} (sarabarāju) — supplies; a supply}

\label{lesson:KA-C470-sarabaraju}



\begin{quote}
Before the new one: say the Kannada for an entrance gate, then the Kannada for a power cut.
\end{quote}

\subsection*{You'll want to know: \kn{ಸರಬರಾಜು}}
\textbf{\kn{ಸರಬರಾಜು}} — \emph{sarabarāju} — \textquotedblleft{}supplies; a supply\textquotedblright{}.

\begin{quote}
\kn{ಮಂಗಳವಾರ} \kn{ನಮ್ಮ} \kn{ಬಡಾವಣೆಗೆ} \kn{ನೀರಿನ} \kn{ಸರಬರಾಜು} \kn{ಇರುವುದಿಲ್ಲ}.
\end{quote}

\begin{quote}
\emph{maṅgaḷavāra namma baḍāvaṇege nīrina sarabarāju iruvudilla.}
\end{quote}

\begin{quote}
On Tuesday there will be no water supply to our area.
\end{quote}

\kn{ಸರಬರಾಜು} \kn{ಮಾಡು}, to supply.

\begin{cousinweb}
\textbf{\kn{ಸರಬರಾಜು}} (\emph{sarabarāju}) is usually traced to Persian \emph{sarbarāh}, \textquotedblleft{}a manager, one who provides\textquotedblright{}, which reached Kannada through Hindustani.
\end{cousinweb}

\begin{grammarlens}[title={What you've built}]
One more word for getting around.
\end{grammarlens}

\subsection*{Your turn}
\begin{itemize}
\raggedright
  \item \emph{Say it:} \emph{sarabarāju}
  \item \emph{Say it:} \emph{sarabarāju}, once more
  \item \emph{Say it:} \emph{maṅgaḷavāra namma baḍāvaṇege nīrina sarabarāju iruvudilla.} — On Tuesday there will be no water supply to our area.
  \item \emph{Recall:} say the Kannada for an entrance gate, then the Kannada for a power cut, then say \emph{sarabarāju} again
\end{itemize}

\begin{tcolorbox}[breakable,colback=teal!4,colframe=teal!35!black,title={Before you move on}]
What does \kn{ಸರಬರಾಜು} mean? (\textquotedblleft{}Supplies; a supply\textquotedblright{}.) Say it once more.
\end{tcolorbox}
\section[\texorpdfstring{\kn{ದೂರವಾಣಿ} (dūravāṇi)}{ದೂರವಾಣಿ (dūravāṇi)}]{\kn{ದೂರವಾಣಿ} (dūravāṇi) — a telephone}

\label{lesson:KA-C470-duravani}



\begin{quote}
Before the new one: say the Kannada for a power cut, then the Kannada for supplies.
\end{quote}

\subsection*{You'll want to know: \kn{ದೂರವಾಣಿ}}
\textbf{\kn{ದೂರವಾಣಿ}} — \emph{dūravāṇi} — \textquotedblleft{}a telephone\textquotedblright{}.

\begin{quote}
\kn{ಕರೆಂಟ್} \kn{ಇಲ್ಲದಿದ್ದರೂ} \kn{ದೂರವಾಣಿ} \kn{ಕೆಲಸ} \kn{ಮಾಡುತ್ತಿತ್ತು}.
\end{quote}

\begin{quote}
\emph{kareṇṭ illadiddarū dūravāṇi kelasa māḍuttittu.}
\end{quote}

\begin{quote}
Even without power the telephone was working.
\end{quote}

\kn{ದೂರ}, far, plus \kn{ವಾಣಿ}, voice: a coined word; in speech people say \kn{ಫೋನ್}.

\begin{cousinweb}
\textbf{\kn{ದೂರವಾಣಿ}} (\emph{dūravāṇi}) is a modern compound built from Sanskrit words: \emph{dūra}, \textquotedblleft{}distant\textquotedblright{} (CDIAL 6495), and \emph{vāṇī}, \textquotedblleft{}voice, sound\textquotedblright{} (CDIAL 11490).
\end{cousinweb}

\begin{grammarlens}[title={What you've built}]
One more word for getting around.
\end{grammarlens}

\subsection*{Your turn}
\begin{itemize}
\raggedright
  \item \emph{Say it:} \emph{dūravāṇi}
  \item \emph{Say it:} \emph{dūravāṇi}, once more
  \item \emph{Say it:} \emph{kareṇṭ illadiddarū dūravāṇi kelasa māḍuttittu.} — Even without power the telephone was working.
  \item \emph{Recall:} say the Kannada for a power cut, then the Kannada for supplies, then say \emph{dūravāṇi} again
\end{itemize}

\begin{tcolorbox}[breakable,colback=teal!4,colframe=teal!35!black,title={Before you move on}]
What does \kn{ದೂರವಾಣಿ} mean? (\textquotedblleft{}A telephone\textquotedblright{}.) Say it once more.
\end{tcolorbox}
